\section{Related Work}
\label{sec:related}

\para{Audience-aware generation and pragmatics.} The Rational Speech Acts tradition models a
speaker who chooses each utterance by simulating how a listener will interpret it relative to
alternatives~\citep{andreas2016reasoning,fried2023pragmatics}. \citet{andreas2016reasoning} sample
candidate descriptions and rerank them with a simulated listener, gaining accuracy while a small
fluency weight keeps language natural; notably, the reranking reasoning did not distill into a
feedforward network, suggesting explicit inference-time reasoning matters. \citet{takmaz2023speaking}
steer a frozen speaker at decode time with a listener simulator, and find that modeling a
\emph{specific} reader (audience-aware) beats modeling a generic one (self-aware), which beats no
modeling at all. \citet{coderesa2025} reverse-generate what an audience thinks a program does and
rerank for distinctiveness, and \citet{crsa2025} maintain a formal per-turn belief over a
listener's private state, gaining most when that belief is genuinely uncertain.
\citet{styleinfusion2023} treat audience preference as an external reward. Across all of these,
audience modeling lives \emph{outside} the generator's own reasoning---as a gradient, a reranker,
a reward, or a belief-update module. Unlike these methods, we place audience modeling
\emph{inside} the generator's natural-language chain-of-thought, and ask whether doing so
per-sentence improves open-ended prose.

\para{Per-sentence and reflective mental-state simulation.} A separate line applies incremental
mental-state reasoning inside the model's own text. \citet{perceptom2024} prompt the model, per
sentence, for who perceived each piece of content, filter to the target's perspective, and then
answer, yielding large false-belief gains. \citet{thoughttracing2025} trace mental states with an
incremental, hypothesis-driven agent, and \citet{tomagent2025} infer a partner's
belief-desire-intention, predict their reply, and reflect. These are the closest method analogs to
our intervention, but all target \emph{comprehension QA or dialogue}, where the goal is to answer
a question about a mind, not to \emph{write well}. We port the per-sentence template to writing and
find that what helps comprehension hurts prose.

\para{Theory of Mind in LLMs, and the chain-of-thought caution.} Default LLM \tom is real but
brittle, collapsing under trivial adversarial perturbations~\citep{ullman2023large,
shapira2023clever}. Crucially, \citet{shapira2023clever} warn that chain-of-thought can inflate
\tom scores through leaked task structure and shortcuts, so a positive result from an
audience-simulation prompt must be checked against controls and a separate judge rather than taken
at face value. Perspective-taking prompts such as \simtom already help \tom
QA~\citep{tomsurvey2025}, and pragmatic competence grows across training
stages~\citep{altprag2025}. We heed the caution directly: we judge only clean extracted text with
judges in different families from the generator, and add an adaptation-gap control that rules out a
generic-quality confound.

\para{Writing quality and homogenization.} LLM prose is fluent but homogenized, and strong judges
struggle to separate good writing from bad~\citep{homogenization2026}. Most relevant to us,
\citet{chakrabarty2025aislop} show that generic chain-of-thought and reasoning models do
\emph{not} improve writing quality, while span-level self-editing does. This makes our question
sharp: if audience-specific reasoning helped writing, it would be a non-trivial exception to the
``more reasoning does not help writing'' finding. Building on their testbed and the homogenization
metrics of \citet{homogenization2026}, we measure quality, audience fit, adaptation, diversity,
and cost together. Our result extends theirs: even audience-specific reasoning fails to help when
applied at fine (per-sentence) granularity, whereas a coarse one-shot audience model does help.
